% =====================================================================
%  Blocos do padrao oficio, montados a partir dos metadados.
%
%  As macros \lkAlgumaCoisa vem de config/metadata.tex, que e gerado.
%  \lkIfAlgumaCoisa{ha valor}{esta vazio} faz cada bloco desaparecer
%  quando o campo correspondente nao foi preenchido — e o que permite
%  usar o mesmo template para um oficio completo e para uma carta
%  simples, sem sobrar rotulo vazio na pagina.
%
%  Para mudar um valor, edite latexkit.config.json.
%  Para mudar a aparencia, use config/style.tex.
% =====================================================================

% --- Timbre ----------------------------------------------------------
% Logotipo (opcional) e identificacao do remetente, centralizados.
\newcommand{\lkTimbre}{%
  \begin{center}
    \lkIfLogo{%
      \includegraphics[height=1.8cm]{\lkLogo}\\[0.6em]
    }{}%
    {\bfseries\MakeUppercase{\lkSender}}%
    \lkIfSenderUnit{\\ \lkSenderUnit}{}%
  \end{center}
  \vspace{1.5em}
}

% --- Identificacao ---------------------------------------------------
% Tipo e numero a esquerda, local e data a direita, na mesma linha.
\newcommand{\lkIdentificacao}{%
  \noindent
  \begin{minipage}[t]{0.5\textwidth}
    \raggedright
    \lkDocumentType\lkIfNumber{ n\textsuperscript{o} \lkNumber}{}
  \end{minipage}%
  \begin{minipage}[t]{0.5\textwidth}
    \raggedleft
    % Data vazia significa "data da compilacao".
    \lkIfCity{\lkCity, }{}\lkIfDate{\lkDate}{\today}
  \end{minipage}

  \vspace{2em}
}

% --- Destinatario ----------------------------------------------------
% Nome, cargo e endereco, alinhados a esquerda.
\newcommand{\lkDestinatario}{%
  \noindent
  \begin{minipage}{\textwidth}
    \raggedright
    \lkIfRecipientTitle{\lkRecipientTitle\\}{}%
    \lkRecipient
    \lkIfRecipientRole{\\ \lkRecipientRole}{}%
    \lkIfRecipientOrg{\\ \lkRecipientOrg}{}%
    \lkIfRecipientAddress{\\ \lkRecipientAddress}{}%
  \end{minipage}

  \vspace{2em}
}

% --- Assunto ---------------------------------------------------------
% O MRPR pede o assunto em negrito, resumido em uma linha.
\newcommand{\lkAssunto}{%
  \lkIfSubject{%
    \noindent\textbf{Assunto:} \lkSubject

    \vspace{2em}
  }{}%
}

% --- Fecho e assinatura ----------------------------------------------
% "Respeitosamente" para autoridade superior, "Atenciosamente" para
% mesma hierarquia ou inferior — a escolha esta em latexkit.config.json.
\newcommand{\lkFecho}{%
  \vspace{2em}
  \noindent \lkClosing,

  \vspace{3em}
  \begin{center}
    \lkSignerName\\
    \lkIfSignerRole{\lkSignerRole}{}%
  \end{center}
}
